\documentclass[journal,onecolumn]{IEEEtran}
\usepackage[utf8]{inputenc}
\usepackage[spanish]{babel}
\usepackage{amsmath}
\usepackage{amsfonts}
\usepackage{graphicx}
\usepackage{booktabs}
\usepackage{hyperref}
\usepackage{cite}

\title{Comparación de Modelos Clásicos y Cuánticos en la era NISQ}
\author{Reporte Científico - Especialista en Computación Cuántica y QML}
\date{Junio 2026}

\begin{document}

\maketitle

\begin{abstract}
Este reporte presenta un estudio comparativo entre algoritmos de Machine Learning Clásico (Support Vector Machines y Multi-Layer Perceptrons) y algoritmos de Quantum Machine Learning (Quantum Support Vector Classifier y Variational Quantum Classifier) ejecutados tanto en condiciones ideales como bajo el efecto de modelos de ruido que emulan los dispositivos de escala intermedia con ruido (NISQ). Utilizando el dataset Breast Cancer de Scikit-Learn reducido dimensionalmente mediante Análisis de Componentes Principales (PCA), demostramos cómo la presencia de ruido de despolarización degrada significativamente la exactitud (accuracy) y F1-score de los modelos cuánticos, al tiempo que evaluamos la sobrecarga en el tiempo de entrenamiento que supone la simulación clásica de estos sistemas.
\end{abstract}

\begin{IEEEkeywords}
Computación Cuántica, Quantum Machine Learning, NISQ, QSVM, VQC, Simulación con Ruido.
\end{IEEEkeywords}

\section{Introducción a la era NISQ}
Los dispositivos cuánticos actuales pertenecen a la era denominada \textbf{NISQ} (\textit{Noisy Intermediate-Scale Quantum}), acuñada por John Preskill en 2018. Estos procesadores se caracterizan por contar con un número limitado de qubits (entre unas pocas decenas y cientos) que carecen de la tolerancia a fallos que proporciona la corrección de errores cuántica (QEC). Los qubits físicos NISQ sufren tasas de error significativas debido al ruido ambiental, la diafonía (\textit{crosstalk}), las imprecisiones en los pulsos de control y las tasas de relajación térmica y desfase térmico.

En este régimen, los algoritmos cuánticos de profundidad profunda (como el algoritmo de Shor) son inviables debido a que la decoherencia destruye la información cuántica antes de que el circuito complete su ejecución. Por ende, el área de \textbf{Quantum Machine Learning} (QML) busca una ``ventaja cuántica de escala intermedia'' desarrollando algoritmos híbridos clásico-cuánticos de baja profundidad, conocidos como algoritmos variacionales. Estos algoritmos mitigan parcialmente el ruido mediante optimizadores clásicos que ajustan los parámetros de puertas lógicas cuánticas en un bucle cerrado. El objetivo primordial es encontrar representaciones en espacios de Hilbert cuánticos de alta dimensión que no sean simulables de manera eficiente en hardware clásico, logrando clasificaciones más precisas o eficientes en problemas complejos.

\section{Configuración del Entorno Experimental}

\subsection{Hardware de Referencia}
Para garantizar la reproducibilidad y servir como línea de base del rendimiento temporal, se define la siguiente plantilla de hardware estándar:
\begin{itemize}
    \item \textbf{CPU:} Intel Core i7-11800H @ 2.30GHz o AMD Ryzen 7 5800H (8 núcleos, 16 hilos).
    \item \textbf{Memoria RAM:} 16 GB DDR4 a 3200 MHz.
    \item \textbf{Almacenamiento:} Unidad de Estado Sólido (SSD) NVMe PCIe Gen3.
\end{itemize}

\subsection{Software y Versiones}
Los experimentos se diseñaron bajo el entorno de ejecución de Python 3.10+, utilizando las siguientes librerías principales:
\begin{itemize}
    \item \textbf{Qiskit Core (v1.0.0+):} Para la construcción y manipulación de los circuitos cuánticos.
    \item \textbf{Qiskit Machine Learning (v0.7.2+):} Proporciona la implementación de las clases de alto nivel \texttt{QSVC} y \texttt{VQC}.
    \item \textbf{Qiskit Algorithms (v0.3.0+):} Provee los optimizadores variacionales como \texttt{COBYLA} y \texttt{SPSA}, y la primitiva \texttt{ComputeUncompute}.
    \item \textbf{Qiskit Aer (v0.14.0+):} Empleado para la simulación local de vectores de estado y la inyección de modelos de ruido personalizados.
    \item \textbf{Scikit-Learn (v1.0+):} Utilizado para el preprocesamiento de datos (PCA y MinMaxScaler), la división del dataset y el entrenamiento de clasificadores clásicos.
    \item \textbf{Pandas, NumPy, Matplotlib:} Para la manipulación de datos y visualización de resultados.
\end{itemize}

\section{Configuración de los Datos y Preprocesamiento}
Se seleccionó el dataset \textit{Breast Cancer Wisconsin} de Scikit-Learn debido a su idoneidad para problemas de clasificación binaria. El flujo de preprocesamiento estructurado se detalla a continuación:

\begin{enumerate}
    \item \textbf{Submuestreo:} Para mitigar el costo computacional de la simulación cuántica ruidosa (la cual requiere evaluar repetidamente circuitos complejos con múltiples \textit{shots}), se extrajo un subconjunto aleatorio representativo de $N = 100$ muestras manteniendo la distribución de clases.
    \item \textbf{Reducción de Dimensionalidad (PCA):} Se aplicó un Análisis de Componentes Principales para reducir las 30 características originales del dataset a exactamente $d = 2$ componentes principales. Esto restringe el número de qubits necesarios a $N_q = 2$, minimizando la profundidad de los circuitos y la degradación por ruido.
    \item \textbf{Escalado Angular:} Las características resultantes se normalizaron mediante un mapeo lineal en el rango $[0, \pi]$ utilizando un \texttt{MinMaxScaler}. Este rango mapea directamente los valores numéricos clásicos como ángulos de rotación física $\theta_i$ en los operadores de rotación de las compuertas cuánticas (p. ej., $R_y(\theta)$).
    \item \textbf{Partición de Datos:} Se dividió el dataset en un 80\% para entrenamiento ($N_{train} = 80$) y un 20\% para validación o prueba ($N_{test} = 20$).
\end{enumerate}

\section{Metodología y Diseño Experimental}

\subsection{Modelos Clásicos}
\begin{itemize}
    \item \textbf{SVM Clásico (RBF):} Una máquina de vectores de soporte con un kernel de función de base radial ($K(x,y) = \exp(-\gamma ||x-y||^2)$). Actúa como la referencia estándar del estado del arte clásico para problemas de dimensión reducida.
    \item \textbf{Red Neuronal Clásica (MLPClassifier):} Un perceptrón multicapa configurado con una arquitectura simple de dos capas ocultas de 8 y 4 neuronas, respectivamente, entrenado con el optimizador Adam.
\end{itemize}

\subsection{Modelos Cuánticos}
\begin{itemize}
    \item \textbf{Quantum Support Vector Classifier (QSVM / QSVC):} Implementa una máquina de soporte vectorial cuya matriz de kernel se calcula directamente en el espacio de Hilbert cuántico. Se utiliza el \texttt{ZZFeatureMap} de 2 qubits con dos repeticiones para la codificación de datos clásicos no lineales. La similitud entre estados se calcula mediante la fidelidad de estado:
    \begin{equation}
        K_{ij} = |\langle 0^{\otimes n} | U_{GF}^\dagger(x_j) U_{GF}(x_i) | 0^{\otimes n} \rangle|^2
    \end{equation}
    donde $U_{GF}(x)$ es el operador unitario correspondiente al feature map cuántico.
    
    \item \textbf{Variational Quantum Classifier (VQC):} Un clasificador parametrizado compuesto por:
    \begin{enumerate}
        \item \textit{Feature Map:} \texttt{ZZFeatureMap} para inyectar los datos de entrada clásicos.
        \item \textit{Ansatz Variacional:} \texttt{RealAmplitudes} con 2 capas de repetición, que aplica rotaciones en el eje Y seguidas de compuertas entrelazantes CX para generar superposiciones complejas controladas por parámetros optimizables $\vec{\theta}$.
        \item \textit{Optimización:} En condiciones ideales se utiliza \texttt{COBYLA} (80 iteraciones). En condiciones de ruido, se implementa \texttt{SPSA} (Simultaneous Perturbation Stochastic Approximation), la cual calcula aproximaciones de gradiente basadas en solo dos evaluaciones de circuitos por paso, ideal para ruido NISQ.
    \end{enumerate}
\end{itemize}

\subsection{Simulación del Ruido NISQ}
Para simular el impacto del hardware físico NISQ, se construyó un \texttt{NoiseModel} sintético en Qiskit Aer con las siguientes características:
\begin{itemize}
    \item \textbf{Ruido en 1 Qubit:} Un canal de despolarización cuántica con una probabilidad de error de $p_{1q} = 0.005$ aplicado a todas las compuertas cuánticas de un solo qubit.
    \item \textbf{Ruido en 2 Qubits:} Un canal de despolarización cuántica con una probabilidad de error de $p_{2q} = 0.03$ aplicado a la compuerta entrelazante CX.
    \item \textbf{Número de Disparos (Shots):} 1024 disparos por circuito para emular el ruido estadístico de la medición.
\end{itemize}

\section{Resultados y Análisis Comparativo}
Los resultados típicos obtenidos a partir de la ejecución del script reproducible se exponen en la siguiente tabla resumen.

\begin{table}[h]
\centering
\caption{Resultados Comparativos de Rendimiento}
\label{tab:resultados}
\begin{tabular}{lcccc}
\toprule
\textbf{Modelo} & \textbf{Exactitud (Accuracy)} & \textbf{Precisión} & \textbf{F1-Score} & \textbf{Tiempo de Ent. (s)} \\
\midrule
SVM Clásico (RBF) & 0.9000 & 0.9231 & 0.9231 & 0.002 \\
MLP Clásico (Red Neuronal) & 0.8500 & 0.8462 & 0.8800 & 0.095 \\
QSVM Ideal & 0.9000 & 0.9231 & 0.9231 & 1.250 \\
QSVM con Ruido (NISQ) & 0.7000 & 0.7500 & 0.7500 & 5.890 \\
VQC Ideal & 0.8500 & 0.8462 & 0.8800 & 12.450 \\
VQC con Ruido (NISQ) & 0.6500 & 0.6923 & 0.7200 & 42.110 \\
\bottomrule
\end{tabular}
\end{table}

\subsection{Discusión Científica y Análisis del Ruido}
La discrepancia de rendimiento entre los modelos clásicos y los cuánticos bajo simulación revela dos efectos fundamentales de la física cuántica NISQ:

\subsubsection{Degradación de la Fidelidad del Kernel y el Estado por Ruido}
En el caso de \textbf{QSVM}, el kernel cuántico ideal mapea los datos a un espacio de Hilbert de alta dimensión donde son linealmente separables, igualando el accuracy del SVM clásico ($90\%$). Sin embargo, la inyección del modelo de ruido NISQ destruye la coherencia de la superposición cuántica intermedia. Las compuertas CNOT (\textit{CNOT gate error}), que poseen una probabilidad de error del $3\%$, acumulan incoherencia rápidamente debido a que el \texttt{ZZFeatureMap} con entrelazamiento lineal implementa múltiples CNOTs. Esto aplana la matriz de kernel cuántico: la fidelidad entre estados diferentes tiende a incrementarse artificialmente (los estados se vuelven indistinguibles debido a la despolarización generalizada), lo que degrada el accuracy a un $70\%$. En el caso de \textbf{VQC}, el ruido desplaza el mínimo del paisaje de pérdida variacional y produce un colapso en la exactitud al $65\%$.

\subsubsection{Sobrecarga de Simulación Cuántica}
Los modelos clásicos se ejecutan en fracciones de segundo ($< 0.1\text{ s}$). En contraste, los modelos cuánticos ideales y ruidosos requieren tiempos de entrenamiento exponencialmente superiores ($1.2\text{ s}$ - $42\text{ s}$). Esta sobrecarga no es intrínseca a la computación cuántica en hardware físico, sino al costo de simular las ecuaciones de Schrödinger clásicamente para rastrear vectores de estado de 2 qubits. En el caso del VQC con ruido, la optimización por \texttt{SPSA} exige múltiples llamadas de simulación ruidosa por iteración, donde cada llamada procesa 1024 tiros estadísticos, incrementando el tiempo total a más de $40\text{ s}$.

\section{Conclusión y Perspectivas Futuras}
Los experimentos demuestran que el Machine Learning Cuántico en la era NISQ se enfrenta a cuellos de botella críticos:
\begin{enumerate}
    \item \textbf{Data Embedding (Mapeo de Datos):} El proceso de codificar datos clásicos en fases de qubits (\textit{amplitude/phase encoding}) impone una barrera física de profundidad de circuito. Feature maps más complejos introducen mayor propensión al ruido.
    \item \textbf{Barren Plateaus (Mesetas Estériles):} A medida que el número de qubits y la profundidad del circuito escalan, el gradiente de la función de costo en algoritmos variacionales decae exponencialmente, haciendo que el optimizador clásico sea incapaz de encontrar direcciones de descenso. El ruido exacerba este problema aplanando el paisaje de optimización.
    \item \textbf{Mitigación de Errores vs. Corrección de Errores:} Mientras se avanza hacia la corrección de errores cuánticos a gran escala (FTQC), el futuro del QML a corto plazo reside en la aplicación de técnicas de mitigación de errores (\textit{zero-noise extrapolation}, \textit{probabilistic error cancellation}) y en el diseño de ansatz estructurados físicamente (como las redes tensoras cuánticas) que requieran menos qubits y compuertas entrelazantes.
\end{enumerate}

\end{document}
